\makeatletter
\def\input@path{{./}{../}{../../}{../../../}{../../../../}{preamble/}{../preamble/}{../../preamble/}{../../../preamble/}{../../../../preamble/}}
\makeatother
\input{preamble_exams.tex}
\input{graphs/G11_Histograms.tex}

\title{Grade 11 -- Term 1 \\ C5, C6 and C7 Exam Solutions -- Version 2}
\author{}
\date{}

\begin{document}
\maketitle
\vspace{-2cm}
\setcounter{tocdepth}{1}
\setcounter{secnumdepth}{2}
\phantomsection\label{toc}
\tableofcontents

\section*{Learning objective}\vspace{-5mm}
Use continuous probability models and sample data to describe typical values,
variability, and interval probabilities, and interpret differences between
sample estimates and theoretical values.

\section*{Criteria assessment}\vspace{-5mm}
This exam assesses C5, C6 and C7 separately, with 6 points per criterion
(18 points in total). Each task shows its point allocation. Partial credit
may be awarded when the correct procedure is clearly shown but arithmetic
errors are made. Answers without the required procedure will not receive credit.
Failure to write your name on the task sheet or submitting it after the deadline
will result in a one-point deduction from each criterion.
For this exam, a minimum of 4 out of 6 points is required to pass each criterion.
Use unrounded intermediate values; round final statistics to three decimal
places and percentage errors to two decimal places where necessary.

\section*{Evaluated criteria}
\begin{enumerate}[label=\textbf{C\arabic*},start=5]
\item Calculates and interprets the mean, variance and standard deviation of continuous distributions using standard formulas and functions.
\item Estimates interval probabilities using relative frequencies from generated samples.
 and compares them with theoretical probabilities.
\item Calculates sample means, variances, and standard deviations for different sample sizes and data representations.
\end{enumerate}


\sectionScored{Two custom-print studios}{C5-\sml{2}}{sec:c5-1}
{
Two studios print custom posters. Each order takes between 3 and 15 hours.
The owners use smooth outlines to describe how often different completion times
occur over many orders; times outside these limits never occur.
At studio I, longer jobs are more common: the outline is a single straight
edge rising from zero at 3 hours to its highest point at 15 hours.
At studio D, shorter jobs are more common: the outline is a single straight
edge falling from its highest point at 3 hours to zero at 15 hours.
Let $I$ and $D$ be the respective completion times in hours.
\QuestionsTasks[sec:c5-1]{
{5}/{sol:c5-1a}/{Identify each distribution from the clues and calculate its mean, variance, and standard deviation using the corresponding formulas. \score{2}},
}
\phantomsection\label{sol:c5-1a}
\subsubsection*{C5.1(a) -- Linear models \score{2}}
The straight rising edge identifies a linearly increasing distribution;
the straight falling edge identifies a linearly decreasing distribution.
Thus $I\sim\operatorname{Lin}_{\mathrm{inc}}(3,15)$ and
$D\sim\operatorname{Lin}_{\mathrm{dec}}(3,15)$, with $a=3$, $b=15$.
\[
E(I)=\frac{a+2b}{3}=\frac{3+2(15)}{3}=11.000\text{ h},\qquad
E(D)=\frac{2a+b}{3}=\frac{2(3)+15}{3}=7.000\text{ h}.
\]
\[
Var(I)=Var(D)=\frac{(b-a)^2}{18}=\frac{12^2}{18}=8.000\text{ h}^2,
\qquad\sigma_I=\sigma_D=\sqrt{8}\approx2.828\text{ h}.
\]

}

\sectionScored{Bookings and daily takings}{C5-\sml{2}}{sec:c5-2}
{
A camera repair shop receives online bookings throughout its opening hours,
averaging 10 bookings per hour. Bookings arrive independently, one at a time,
at a steady pace. In a very short time window, the chance of a booking is
proportional to the window's length; two or more bookings are negligible.
Past waiting does not affect the chance of a booking during the next minute.
Let $T$ be the time between consecutive bookings in minutes.

On comparable working days, total revenue varies because of many small,
independent influences, with none dominating. Its smooth outline has a symmetric
bell shape, highest at \$3200 and tapering on both sides. Its sides change from
becoming steeper to becoming flatter at \$3000 and \$3400.
Treat this bell-shaped pattern as exact. Let $R$ be daily revenue in hundreds
of dollars.
\QuestionsTasks[sec:c5-2]{
{5}/{sol:c5-2a}/{Identify the distribution for T from the clues and use its formulas to find the mean, variance, and standard deviation in minutes. \score{1}},
{5}/{sol:c5-2b}/{Identify the revenue distribution from the clues. Use its formulas to find the mean, variance, and standard deviation in dollars. Interpret the mean. \score{1}}
}
\phantomsection\label{sol:c5-2a}
\subsubsection*{C5.2(a) -- Exponential waiting time \score{1}}
Independent single arrivals at a steady pace and memoryless waiting identify
an exponential distribution. The rate is $\lambda=10/60=1/6$ per minute,
so $T\sim\operatorname{Exp}(1/6)$.
\[
E(T)=\frac{1}{\lambda}=6.000\text{ min},\qquad
Var(T)=\frac{1}{\lambda^2}=36.000\text{ min}^2,\qquad
\sigma_T=\frac{1}{\lambda}=6.000\text{ min}.
\]
\phantomsection\label{sol:c5-2b}
\subsubsection*{C5.2(b) -- Normal revenue and units \score{1}}
The symmetric bell shape and independent small influences identify a normal
distribution. Its mean is \$3200. Curvature changes occur one standard deviation
from the mean: $\sigma_Y=3400-3200=3200-3000=200$ dollars.
In hundreds of dollars, $\mu_R=32$, $\sigma_R=2$, and $\sigma_R^2=4$;
thus $R\sim N(32,4)$ (mean, variance). Let $Y=100R$ be revenue in dollars.
\[
E(Y)=100\mu_R=100(32)=3200.000\text{ dollars},\qquad
Var(Y)=100^2\sigma_R^2=100^2(4)=40000.000\text{ dollars}^2,
\]
\[
\sigma_Y=\sqrt{Var(Y)}=\sqrt{40000}=200.000\text{ dollars}.
\]
The long-run average daily revenue is \$3200; individual days may differ.
\secInicio[sec:c5-2]

}

\sectionScored{Two parcel-labeling stations}{C5-\sml{2}}{sec:c5-3}
{
A delivery company uses two parcel-labeling stations. At station A, an automated
machine labels a parcel in anywhere from 5 to 19 minutes. Within these limits,
any two time windows of the same length are equally likely to contain the
finishing time. Let $A$ be its labeling time in minutes.

At station B, staff label special parcels, also taking between 5 and 19 minutes.
The manager draws a smooth outline of how often different labeling times occur.
The outline starts at zero at 5 minutes, rises along a straight edge to its
single highest point at 10 minutes, then falls along a straight edge to zero
at 19 minutes. No labeling times occur outside these limits.
Let $B$ be this labeling time in minutes.
\QuestionsTasks[sec:c5-3]{
{5}/{sol:c5-3a}/{Identify the distribution for A from the clues and calculate its mean, variance, and standard deviation using the corresponding formulas. \score{1}},
{5}/{sol:c5-3b}/{Identify the distribution for B from the clues and calculate its mean, variance, and standard deviation using the corresponding formulas. \score{1}}
}
\phantomsection\label{sol:c5-3a}
\subsubsection*{C5.3(a) -- Uniform model \score{1}}
Equal-length windows being equally likely identifies $A\sim U(5,19)$.
With $a=5$, $b=19$, apply $\mu=(a+b)/2$, $\sigma^2=(b-a)^2/12$,
and $\sigma=\sqrt{\sigma^2}$:
\[
E(A)=\frac{5+19}{2}=12.000\text{ min},\qquad
Var(A)=\frac{(19-5)^2}{12}=\frac{49}{3}\approx16.333\text{ min}^2,
\qquad\sigma_A=\sqrt{49/3}\approx4.041\text{ min}.
\]
\phantomsection\label{sol:c5-3b}
\subsubsection*{C5.3(b) -- Triangular model \score{1}}
Two straight edges meeting at one peak identify $B\sim\operatorname{Tri}(5,19,10)$
(minimum, maximum, mode). Apply $\mu=(a+b+c)/3$ and
$\sigma^2=(a^2+b^2+c^2-ab-ac-bc)/18$:
\[
E(B)=\frac{5+19+10}{3}=\frac{34}{3}\approx11.333\text{ min},
\]
\[
Var(B)=\frac{5^2+19^2+10^2-5(19)-5(10)-19(10)}{18}
=\frac{486-335}{18}=\frac{151}{18}\approx8.389\text{ min}^2,
\qquad\sigma_B=\sqrt{151/18}\approx2.896\text{ min}.
\]

}

\sectionScored{File upload times, grouped data}{C6-\sml{2}}{sec:c6-1}
{
A sample of file upload times in seconds is summarized below.
Classes are left-inclusive and right-exclusive.
\[
\begin{array}{c|rrrr}
\text{Class}&[2,8)&[8,14)&[14,20)&[20,26)\\\hline
f_i&5&9&7&3
\end{array}
\]
The theoretical probability is $P(8\le X<20)=0.75$.
\QuestionsTasks[sec:c6-1]{
{6}/{sol:c6-1a}/{Find the empirical probability of the stated interval from the given sample. \score{1}},
{6}/{sol:c6-1b}/{Calculate the absolute and percentage errors relative to the given theoretical probability. \score{1}}
}
\phantomsection\label{sol:c6-1a}
\subsubsection*{C6.1(a) -- Relative frequency \score{1}}
The event contains the two middle classes:
$n=5+9+7+3=24$ and $k=9+7=16$.
Because the event is a union of complete classes, no individual upload
times need to be estimated.
\[
\widehat P(8\le X<20)=\frac{16}{24}=\frac{2}{3}\approx0.667.
\]
\phantomsection\label{sol:c6-1b}
\subsubsection*{C6.1(b) -- Errors \score{1}}
Using the exact relative frequency,
\[
E=\left|\frac{2}{3}-0.75\right|=\frac{1}{12}\approx0.083,\qquad
\text{Percentage error}=\frac{\frac{1}{12}}{0.75}\times100\%\approx11.11\%.
\]
\secInicio[sec:c6-1]

}

\sectionScored{Office electricity use, raw data}{C6-\sml{1}}{sec:c6-2}
{
A simulated sample of daily office electricity use (tens of kilowatt-hours) is
\[
\begin{array}{cccccccccc}
1.3&1.7&2.0&2.2&2.6&2.9&3.3&3.5&3.8&4.2\\
1.6&1.9&2.3&2.7&3.1&3.4&3.6&3.9&4.4&4.6
\end{array}
\]
The theoretical probability is $P(2\le X\le3.5)=0.60$.
\QuestionsTasks[sec:c6-2]{
{6}/{sol:c6-2a}/{Find the empirical probability of the stated interval from the given sample. \score{0.5}},
{6}/{sol:c6-2b}/{Calculate the absolute and percentage errors relative to the given theoretical probability. \score{0.5}}
}
\phantomsection\label{sol:c6-2a}
\subsubsection*{C6.2(a) -- Relative frequency \score{0.5}}
There are $n=20$ observations. The included observations are
$2.0,2.2,2.6,2.9,3.3,3.5,2.3,2.7,3.1,3.4$, so $k=10$.
\[
\widehat P(2\le X\le3.5)=\frac{10}{20}=\frac{1}{2}\approx0.500.
\]
\phantomsection\label{sol:c6-2b}
\subsubsection*{C6.2(b) -- Errors \score{0.5}}
Using the exact relative frequency,
\[
E=\left|\frac{1}{2}-0.60\right|=\frac{1}{10}\approx0.100,\qquad
\text{Percentage error}=\frac{\frac{1}{10}}{0.60}\times100\%\approx16.67\%.
\]
\secInicio[sec:c6-2]

}

\sectionScored{Museum entrance queues, histogram data}{C6-\sml{2}}{sec:c6-3}
{
The histogram is the only record of a sample of museum entrance queue times.
Bar heights are frequencies; times are in minutes, with left-inclusive,
right-exclusive classes.
\begin{center}
\HistogramGraph{(1,5),(5,9),(9,13),(13,17)}{8,15,11,6}
{0.60\textwidth}{4.5cm}{blue!25}{4}{5}
\end{center}
The theoretical probability is $P(5\le X<13)=0.70$.
\QuestionsTasks[sec:c6-3]{
{6}/{sol:c6-3a}/{Find the empirical probability of the stated interval from the given sample. \score{1}},
{6}/{sol:c6-3b}/{Calculate the absolute and percentage errors relative to the given theoretical probability. \score{1}}
}
\phantomsection\label{sol:c6-3a}
\subsubsection*{C6.3(a) -- Relative frequency \score{1}}
The classes $[1,5),[5,9),[9,13),[13,17)$ have frequencies
$8,15,11,6$. Thus $n=8+15+11+6=40$ and $k=15+11=26$.
\[
\widehat P(5\le X<13)=\frac{26}{40}=\frac{13}{20}\approx0.650.
\]
\phantomsection\label{sol:c6-3b}
\subsubsection*{C6.3(b) -- Errors \score{1}}
Using the exact relative frequency,
\[
E=\left|\frac{13}{20}-0.70\right|=\frac{1}{20}\approx0.050,\qquad
\text{Percentage error}=\frac{\frac{1}{20}}{0.70}\times100\%\approx7.14\%.
\]
\secInicio[sec:c6-3]

}

\sectionScored{Equipment rental times, frequency data}{C6-\sml{1}}{sec:c6-4}
{
Recorded equipment rental times are summarized below.
\[
\begin{array}{c|rrrrrr}
x_i\text{ (min)}&4&8&12&16&20&24\\\hline
f_i&8&16&28&22&14&12
\end{array}
\]
The theoretical probability is $P(8\le X<20)=0.80$.
\QuestionsTasks[sec:c6-4]{
{6}/{sol:c6-4a}/{Find the empirical probability of the stated interval from the given sample. \score{0.5}},
{6}/{sol:c6-4b}/{Calculate the absolute and percentage errors relative to the given theoretical probability. \score{0.5}}
}
\phantomsection\label{sol:c6-4a}
\subsubsection*{C6.4(a) -- Relative frequency \score{0.5}}
The total is $n=8+16+28+22+14+12=100$.
The value 20 is excluded, so $k=16+28+22=66$.
\[
\widehat P(8\le X<20)=\frac{66}{100}=\frac{33}{50}\approx0.660.
\]
\phantomsection\label{sol:c6-4b}
\subsubsection*{C6.4(b) -- Errors \score{0.5}}
Using the exact relative frequency,
\[
E=\left|\frac{33}{50}-0.80\right|=\frac{7}{50}\approx0.140,\qquad
\text{Percentage error}=\frac{\frac{7}{50}}{0.80}\times100\%\approx17.50\%.
\]
\secInicio[sec:c6-4]

}

Use the sample-variance denominator $n-1$ throughout these four problems.
For each statistic, calculate absolute error $E=|\text{sample value}-\text{theoretical value}|$
and percentage error $100E/|\text{theoretical value}|$. Use unrounded values in all error calculations.
For grouped data and histograms, use class midpoints to estimate the statistics.

\sectionScored{Luggage masses, histogram data}{C7-\sml{2}}{sec:c7-1}
{
A histogram gives a sample of luggage masses in kilograms. Bar heights are
frequencies; classes are left-inclusive and right-exclusive.
The model has $\mu=9$ kg and $\sigma^2=7$ kg$^2$.
\begin{center}
\HistogramGraph{(4,7),(7,10),(10,13)}{3,4,2}
{0.55\textwidth}{4.5cm}{orange!30}{3}{2}
\end{center}
\QuestionsTasks[sec:c7-1]{
{7}/{sol:c7-1a}/{Calculate the sample mean, sample variance, and sample standard deviation from the given data. \score{1}},
{7}/{sol:c7-1b}/{Calculate the absolute and percentage errors of the mean, variance, and standard deviation relative to the theoretical values. \score{1}}
}
\phantomsection\label{sol:c7-1a}
\subsubsection*{C7.1(a) -- Histogram statistics \score{1}}
The class midpoints are $5.5,8.5,11.5$ kg. Use these as the representative observations.
\[
n=3+4+2=9,\qquad
\bar{x}\approx\frac{3(\frac{11}{2})+4(\frac{17}{2})+2(\frac{23}{2})}{9}=\frac{49}{6}\approx8.167.
\]
The weighted squared-deviation sum, using the unrounded mean, is
\[
3\left(\frac{11}{2}-\frac{49}{6}\right)^2+4\left(\frac{17}{2}-\frac{49}{6}\right)^2+2\left(\frac{23}{2}-\frac{49}{6}\right)^2=44.
\]
\[
s^2\approx\frac{44}{9-1}=\frac{11}{2}\approx5.500,\qquad
s\approx\sqrt{\frac{11}{2}}\approx2.345.
\]
The mean and standard deviation are in kg; the variance is in squared kg.

\phantomsection\label{sol:c7-1b}
\subsubsection*{C7.1(b) -- Absolute and percentage errors \score{1}}
The theoretical standard deviation is $\sigma=\sqrt{7}$.
Apply $E=|\text{sample value}-\text{theoretical value}|$ and
$100E/|\text{theoretical value}|$ to the unrounded results:
\[
\begin{array}{l|rr}
\text{Statistic}&\text{Absolute error}&\text{Percentage error}\\\hline
Mean&0.833&9.26\%\\
Variance&1.500&21.43\%\\
Standard deviation&0.301&11.36\%\\
\end{array}
\]
Absolute errors have the units of the corresponding statistic; percentage errors are unitless.
These errors use midpoint estimates because individual observations are unknown.
\secInicio[sec:c7-1]

}

\sectionScored{Chamber temperatures, raw data}{C7-\sml{1}}{sec:c7-2}
{
Five chamber temperature readings in degrees Celsius are $14,17,18,19,22$.
The theoretical model has $\mu=19$ degrees Celsius and $\sigma^2=6$ square degrees Celsius.
\QuestionsTasks[sec:c7-2]{
{7}/{sol:c7-2a}/{Calculate the sample mean, sample variance, and sample standard deviation from the given data. \score{0.5}},
{7}/{sol:c7-2b}/{Calculate the absolute and percentage errors of the mean, variance, and standard deviation relative to the theoretical values. \score{0.5}}
}
\phantomsection\label{sol:c7-2a}
\subsubsection*{C7.2(a) -- Raw-data statistics \score{0.5}}
\[
n=1+1+1+1+1=5,\qquad
\bar{x}=\frac{1(14)+1(17)+1(18)+1(19)+1(22)}{5}=18\approx18.000.
\]
\[
\begin{array}{c|rrrrr}
x_i&14&17&18&19&22\\\hline
x_i-\bar{x}&-4&-1&0&1&4\\
(x_i-\bar{x})^2&16&1&0&1&16
\end{array}
\]
The squared deviations sum to $16+1+0+1+16=34$.
\[
s^2=\frac{34}{5-1}=\frac{17}{2}\approx8.500,\qquad
s=\sqrt{\frac{17}{2}}\approx2.915.
\]
The mean and standard deviation are in degrees Celsius; the variance is in squared degrees Celsius.

\phantomsection\label{sol:c7-2b}
\subsubsection*{C7.2(b) -- Absolute and percentage errors \score{0.5}}
The theoretical standard deviation is $\sigma=\sqrt{6}$.
Apply $E=|\text{sample value}-\text{theoretical value}|$ and
$100E/|\text{theoretical value}|$ to the unrounded results:
\[
\begin{array}{l|rr}
\text{Statistic}&\text{Absolute error}&\text{Percentage error}\\\hline
Mean&1.000&5.26\%\\
Variance&2.500&41.67\%\\
Standard deviation&0.466&19.02\%\\
\end{array}
\]
Absolute errors have the units of the corresponding statistic; percentage errors are unitless.
\secInicio[sec:c7-2]

}

\sectionScored{Parcels dispatched, frequency data}{C7-\sml{1}}{sec:c7-3}
{
The numbers of parcels dispatched in nine short collection periods are summarized below.
The theoretical model has $\mu=7$ parcels and $\sigma^2=9$ parcels$^2$.
\[
\begin{array}{c|rrr}x_i&3&7&11\\\hline f_i&2&4&3\end{array}
\]
\QuestionsTasks[sec:c7-3]{
{7}/{sol:c7-3a}/{Calculate the sample mean, sample variance, and sample standard deviation from the given data. \score{0.5}},
{7}/{sol:c7-3b}/{Calculate the absolute and percentage errors of the mean, variance, and standard deviation relative to the theoretical values. \score{0.5}}
}
\phantomsection\label{sol:c7-3a}
\subsubsection*{C7.3(a) -- Weighted statistics \score{0.5}}
\[
n=2+4+3=9,\qquad
\bar{x}=\frac{2(3)+4(7)+3(11)}{9}=\frac{67}{9}\approx7.444.
\]
The weighted squared-deviation sum, using the unrounded mean, is
\[
2\left(3-\frac{67}{9}\right)^2+4\left(7-\frac{67}{9}\right)^2+3\left(11-\frac{67}{9}\right)^2=\frac{704}{9}.
\]
\[
s^2=\frac{\frac{704}{9}}{9-1}=\frac{88}{9}\approx9.778,\qquad
s=\sqrt{\frac{88}{9}}\approx3.127.
\]
The mean and standard deviation are in parcels; the variance is in squared parcels.

\phantomsection\label{sol:c7-3b}
\subsubsection*{C7.3(b) -- Absolute and percentage errors \score{0.5}}
The theoretical standard deviation is $\sigma=\sqrt{9}$.
Apply $E=|\text{sample value}-\text{theoretical value}|$ and
$100E/|\text{theoretical value}|$ to the unrounded results:
\[
\begin{array}{l|rr}
\text{Statistic}&\text{Absolute error}&\text{Percentage error}\\\hline
Mean&0.444&6.35\%\\
Variance&0.778&8.64\%\\
Standard deviation&0.127&4.23\%\\
\end{array}
\]
Absolute errors have the units of the corresponding statistic; percentage errors are unitless.
\secInicio[sec:c7-3]

}

\sectionScored{Walking times, grouped data}{C7-\sml{2}}{sec:c7-4}
{
Walking times in minutes are grouped using left-inclusive, right-exclusive classes.
The theoretical mean is 23 minutes and variance is 30 minutes$^2$.
\[
\begin{array}{c|rrr}
\text{Time (min)}&[12,18)&[18,24)&[24,30)\\\hline
f_i&2&3&4
\end{array}
\]
\QuestionsTasks[sec:c7-4]{
{7}/{sol:c7-4a}/{Calculate the sample mean, sample variance, and sample standard deviation from the given data. \score{1}},
{7}/{sol:c7-4b}/{Calculate the absolute and percentage errors of the mean, variance, and standard deviation relative to the theoretical values. \score{1}}
}
\phantomsection\label{sol:c7-4a}
\subsubsection*{C7.4(a) -- Midpoint estimates \score{1}}
The class midpoints are $15,21,27$ min. Use these as the representative observations.
\[
n=2+3+4=9,\qquad
\bar{x}\approx\frac{2(15)+3(21)+4(27)}{9}=\frac{67}{3}\approx22.333.
\]
The weighted squared-deviation sum, using the unrounded mean, is
\[
2\left(15-\frac{67}{3}\right)^2+3\left(21-\frac{67}{3}\right)^2+4\left(27-\frac{67}{3}\right)^2=200.
\]
\[
s^2\approx\frac{200}{9-1}=25\approx25.000,\qquad
s\approx\sqrt{25}\approx5.000.
\]
The mean and standard deviation are in min; the variance is in squared min.

\phantomsection\label{sol:c7-4b}
\subsubsection*{C7.4(b) -- Absolute and percentage errors \score{1}}
The theoretical standard deviation is $\sigma=\sqrt{30}$.
Apply $E=|\text{sample value}-\text{theoretical value}|$ and
$100E/|\text{theoretical value}|$ to the unrounded results:
\[
\begin{array}{l|rr}
\text{Statistic}&\text{Absolute error}&\text{Percentage error}\\\hline
Mean&0.667&2.90\%\\
Variance&5.000&16.67\%\\
Standard deviation&0.477&8.71\%\\
\end{array}
\]
Absolute errors have the units of the corresponding statistic; percentage errors are unitless.
These errors use midpoint estimates because individual observations are unknown.
\secInicio[sec:c7-4]


\newpage
\section{Summarized Solutions}
\subsection*{C5}
\begin{enumerate}[label=\textbf{C5.\arabic*.}]
\item (a) I is linearly increasing and D linearly decreasing on $[3,15]$.
Means: $11.000$ h and $7.000$ h; both variances $8.000$ h$^2$ and SDs $2.828$ h.
\item (a) $T\sim\operatorname{Exp}(1/6)$ in minutes; mean $6.000$ min,
variance $36.000$ min$^2$, SD $6.000$ min.
(b) $R\sim N(32,4)$ (mean, variance) in hundreds of dollars.
Dollar revenue: mean $3200.000$, variance $40000.000$, SD $200.000$;
long-run daily average \$3200.
\item (a) $A\sim U(5,19)$; mean $12.000$ min, variance $16.333$ min$^2$, SD $4.041$ min.
(b) $B\sim\operatorname{Tri}(5,19,10)$ (minimum, maximum, mode);
mean $11.333$ min, variance $8.389$ min$^2$, SD $2.896$ min.
\end{enumerate}
\subsection*{C6}
\begin{enumerate}[label=\textbf{C6.\arabic*.}]
\item (a) $n=24$, $k=16$, $\widehat P\approx0.667$.
(b) Absolute error $\approx0.083$; percentage error $\approx11.11\%$.
\item (a) $n=20$, $k=10$, $\widehat P\approx0.500$.
(b) Absolute error $\approx0.100$; percentage error $\approx16.67\%$.
\item (a) $n=40$, $k=26$, $\widehat P\approx0.650$.
(b) Absolute error $\approx0.050$; percentage error $\approx7.14\%$.
\item (a) $n=100$, $k=66$, $\widehat P\approx0.660$.
(b) Absolute error $\approx0.140$; percentage error $\approx17.50\%$.
\end{enumerate}
\subsection*{C7}
\begin{enumerate}[label=\textbf{C7.\arabic*.}]
\item (a) $n=9$, $\bar{x}\approx8.167$, $s^2\approx5.500$,
$s\approx2.345$ (mean and SD in kg; variance in squared kg).
Midpoints $(5.5,8.5,11.5)$; statistics are midpoint estimates.
(b) Absolute errors (mean, variance, SD): $(0.833,1.500,0.301)$; percentage errors: $(9.26\%,21.43\%,11.36\%)$.
\item (a) $n=5$, $\bar{x}\approx18.000$, $s^2\approx8.500$,
$s\approx2.915$ (mean and SD in degrees Celsius; variance in squared degrees Celsius).
(b) Absolute errors (mean, variance, SD): $(1.000,2.500,0.466)$; percentage errors: $(5.26\%,41.67\%,19.02\%)$.
\item (a) $n=9$, $\bar{x}\approx7.444$, $s^2\approx9.778$,
$s\approx3.127$ (mean and SD in parcels; variance in squared parcels).
(b) Absolute errors (mean, variance, SD): $(0.444,0.778,0.127)$; percentage errors: $(6.35\%,8.64\%,4.23\%)$.
\item (a) $n=9$, $\bar{x}\approx22.333$, $s^2\approx25.000$,
$s\approx5.000$ (mean and SD in min; variance in squared min).
Midpoints $(15,21,27)$; statistics are midpoint estimates.
(b) Absolute errors (mean, variance, SD): $(0.667,5.000,0.477)$; percentage errors: $(2.90\%,16.67\%,8.71\%)$.
\end{enumerate}
}

\end{document}
